%% Chapter1_Introduction.tex
\chapter{Introduction}
\label{chap:intro}

\section{Background and Motivation}
Many junctions, especially in smaller cities, still run fixed-time signal plans in which each approach receives a preset green interval regardless of how many vehicles are waiting. Such plans are dimensioned for average flows; when demand fluctuates or is uneven, green time is wasted on nearly empty approaches while queues build on busy ones. Traffic in Indian cities adds variety: motorcycles, auto-rickshaws, buses and trucks share the road with cars, and emergency vehicles must cross a junction without delay.

Traffic-responsive systems such as SCATS and SCOOT adjust timings from detector data but rely on dedicated detector infrastructure~\cite{ref3,ref4}. Low-cost alternatives estimate density locally with infrared sensors, image-based density classes or camera detectors~\cite{ref1,ref5,ref6,ref8,ref9,ref10}. These systems are usually demonstrated on a prototype or a few scenarios and report an average improvement over a fixed plan. Four questions are rarely answered: whether variable greens keep safe change intervals; how emergency priority interacts with density-based timing; how the controller behaves as demand approaches and exceeds capacity; and how much of the gain is due to the density logic rather than to replacing an untuned plan with any demand-responsive one.

\section{Problem Statement}
The lack of a rigorous, reproducible evaluation of low-cost density-based signal control, against strong baselines, near and above capacity, and with verified signal safety, leaves it unclear when such a controller is actually better than a well-tuned fixed or actuated plan. This project addresses it by building a complete simulation-first traffic management system and evaluating its density-aware controller on identical seeded traffic against re-implementations of the strategies used by existing approaches.

\section{Research Objectives}
The key objectives of this project are:
\begin{enumerate}
  \item To design and develop a density-aware adaptive signal controller, an emergency-vehicle preemption layer that keeps full change intervals, and a vision pipeline, with a live dashboard and a hardware bridge.
  \item To analyse and evaluate the controller against a fixed plan and four re-implemented strategies (Webster-optimized, actuated, max-pressure and three-level density control) below, at and above capacity.
  \item To validate the system: signal safety at every step, simulated queue discharge against Highway Capacity Manual values, the vision pipeline against exact ground truth, and the effect of imperfect sensing.
  \item To demonstrate the societal relevance of the work in line with the UN Sustainable Development Goals.
\end{enumerate}

\section{Novelty and Originality}
Rather than a new detector, this project studies the \emph{timing logic} of a density-based controller under paired, seeded conditions against re-implementations of a Webster-optimized plan~\cite{ref2}, vehicle-actuated control~\cite{ref23}, max-pressure control~\cite{ref11} and the three-level density control of~\cite{ref1}. It exposes a failure mode of density-only switching near capacity and corrects it with a discharge-hold rule taken from actuated control. The role of AI is stated plainly: the controller is a transparent rule, not a learned policy; machine learning enters through the pre-trained Haar-cascade detector, whose accuracy is measured, and through a data-driven calibration of the score on held-out seeds. The contributions are listed in Section~\ref{sec:contrib}.

\section{Scope of the Project}
\textbf{Included:} a microscopic simulator of a four-arm junction with seven vehicle classes; fixed-time, adaptive and four comparison controllers; emergency preemption with a per-step safety audit; a Haar-cascade detector and ByteTrack-style tracker evaluated on synthetic frames; a FastAPI backend with SQLite persistence; a React dashboard; a Raspberry Pi/ESP32 driver; Docker packaging and CI; and the experiments of Chapter~\ref{chap:results}.

\textbf{Excluded:} turning movements, lane changes, pedestrians and multi-junction coordination; real camera footage; closing the control loop on detector output; reinforcement learning; and tests on a physical junction or board (Section~\ref{sec:future}).

\section{Report Organisation}
Chapter 2 presents the literature review.
Chapter 3 describes the research methodology.
Chapter 4 covers the implementation details.
Chapter 5 presents results and discussion.
Chapter 6 discusses societal impact and SDG alignment.
Chapter 7 addresses Universal Human Values.
Chapter 8 concludes the report with future directions.
